The judicial channel for redistricting reform is uniquely powerful: it does not require incumbent legislators to cooperate and is available in every state where a constitutional violation can be established. Its principal limitations---slow remediation, high cost, and expert-dependent verification---are directly addressed by the \texttt{bisect} platform's algorithmic approach.

The three model court orders in Section~3 provide courts with ready-to-adapt language for each major remediation scenario. Model Order 2 (special-master-input) is the recommended starting point: it uses the algorithm's speed and verifiability to reduce remediation timelines from 14 months to 60 days, while preserving judicial oversight and VRA modification opportunities.

The SHA-256 audit chain directly answers the reproducibility challenge that courts have implicitly imposed on redistricting experts since \textit{LWV v. PA}. Any party can run the verification command; disputes about map authenticity are resolved mechanically. This reduces litigation cost and enables courts to focus on the substantive legal questions---VRA compliance, criteria interpretation, constitutional standards---rather than on competing cartographic expert battles.

The practical recommendation for litigants in active redistricting cases is direct: in any case where the court is likely to order a remedial map, plaintiffs should introduce evidence of the \texttt{bisect} platform early in the proceeding, demonstrate the algorithmic baseline for the state in question, and propose the court adopt Model Order 2 as the remediation framework. This positions the platform as a verified, neutral resource available to the court, not as a partisan tool introduced late in proceedings.

Redistricting litigation is expensive and slow not because the legal questions are novel but because the technical questions (is this map neutral? does it satisfy these criteria?) are resolved through expert judgment that is expensive and contestable. The \texttt{bisect} platform's algorithmic approach makes these technical questions mechanically resolvable, freeing courts and litigants to focus on the law.
